\documentclass[aps,prl,twocolumn,showpacs,superscriptaddress]{revtex4-2}
\usepackage{amsmath,amssymb,booktabs,xcolor,hyperref}

\begin{document}

\begin{abstract}
We derive atomic structure, the Pauli exclusion principle, the octet 
rule, and the chemical inertness of noble gases from BCT vacuum topology, 
without additional postulates. Electrons are H=1 Hopf defects of the 
Octet-Hopfion Condensate (OHC); protons are D4 instanton nodes; electron 
shells are quantised OHC Bessel resonance modes with capacity $2n^2$. 
The Pauli exclusion principle is the topological complementarity of 
$H=+1$ and $H=-1$ Hopfions, which combine to net $H=0$: a theorem, 
not a postulate. Noble gases are \textbf{closed Hopf multiplets} --- 
topologically inert, not merely energetically stable. The octet rule 
derives from the two BCT void types: $r_{\rm oct}$ generates 1 spherical 
(s) mode; $r_{\rm tet}$ generates 3 directional (p) modes; giving 
$4 \times 2 = 8$ electrons at $n=2$. Gilbert Lewis drew $r_{\rm oct}$ 
and $r_{\rm tet}$ in 1916. He called them dots. Chemistry is topology. 
The periodic table is a Hopf fibration atlas.
\end{abstract}

\maketitle

\section{Introduction}

Standard quantum mechanics derives atomic structure correctly but 
incompletely. It provides rules without reasons: the Pauli exclusion 
principle is postulated, not derived; the octet rule is a pattern, 
not a theorem; noble gas inertness is energetic, not topological. 
The BCT Superfluid Lattice Model provides the missing depth.

The recent result of White et al.~\cite{White2026} --- deriving the 
hydrogen spectrum acoustically from a dynamic vacuum via the 
Madelung/Bogoliubov framework --- is the limiting case of BCT. 
BCT specifies the exact vacuum geometry ($c/a = \sqrt{2}$, the unique 
Lorentz-invariant BCT lattice) and derives all coupling constants.

\section{BCT Atomic Constituents}

\begin{center}
\begin{tabular}{lll}
\toprule
Entity & BCT identity & Origin \\
\midrule
Proton/neutron & D4 instanton node & Letter~71 \\
Electron & $H=1$ Hopfion & $\pi_3(S^2)=\mathbb{Z}$ \\
Electron spin & OHC phase $\pm$ & Opposite orientation \\
Shell $n$ & OHC Bessel mode $n$ & Josephson at $\alpha_0$ \\
\bottomrule
\end{tabular}
\end{center}

\section{Hydrogen}

The hydrogen atom: single D4 node (proton) with single $H=1$ Hopfion 
(electron) in the first OHC Bessel resonance. The Bohr radius:
\begin{equation}
    a_0 = \frac{\hbar}{m_e c\,\alpha} = 0.5293~\text{\AA}
    \quad \text{(observed: } 0.5292~\text{\AA}\ \checkmark\text{)}
\end{equation}
where $\alpha = \alpha_0(1-2\alpha_0) = 1/137.018$ (Letter~1, error 
$-0.013\%$). Shell energies $E_n = -13.606/n^2~\text{eV}$ follow from 
the OHC Bessel mode spectrum; the $1/n^2$ ladder is exact.

\section{Pauli Exclusion as Topological Complementarity}

\begin{theorem}[BCT Pauli Theorem]
Two $H=+1$ Hopfions in the same spatial OHC resonance mode must have 
opposite OHC interior phase orientations (opposite spin). Two maps 
$S^3\to S^2$ with the \emph{same} orientation combine to $H=+2$, 
exceeding the single-mode capacity. Two maps with \emph{opposite} 
orientation ($H=+1$ and $H=-1$) combine to net $H=0$ --- topologically 
complementary --- and share the same mode.
\end{theorem}

The Pauli principle is not postulated. It is topological necessity: 
$H=+2$ exceeds the $n=1$ mode capacity. $H\in\mathbb{Z}$ follows from 
$\pi_3(S^2) = \mathbb{Z}$ (Appendix~JH).

\section{The Octet Rule}

The BCT vacuum has exactly two void types. Their geometry directly 
determines the orbital symmetries at $n=2$:
\begin{itemize}
\item $r_{\rm oct}$: 1 isotropic centre $\to$ 1 s-orbital $\to$ 2 electrons
\item $r_{\rm tet}$: 3 axial orientations $\to$ 3 p-orbitals $\to$ 6 electrons
\end{itemize}
Total: $4 \times 2 = \mathbf{8}$ electrons. The octet is the BCT void 
geometry at $n=2$. Lewis drew $r_{\rm oct} + 3r_{\rm tet}$ in 1916. 
He called them dots.

\section{Noble Gases as Closed Hopf Multiplets}

\begin{theorem}[BCT Noble Gas Theorem]
A noble gas is an atom whose outermost electron shell is a \textbf{closed 
Hopf multiplet}: all $2n^2$ Hopf modes at level $n$ are occupied, 
with net shell Hopf charge $= 0$.

A closed Hopf multiplet is \textbf{topologically inert}: it cannot 
accept an additional electron (which would require opening a new Hopf 
mode at the next resonance level) and cannot donate an electron 
(which would break the topological closure) without supplying 
energy $\propto \Delta E_{\rm top}$.
\end{theorem}

Noble gas inertness is not merely energetic. It is topological. 
The Hopf multiplet is complete. There is no available topological 
channel for bonding.

\begin{center}
\begin{tabular}{lllr}
\toprule
Gas & $Z$ & Outer config & BCT status \\
\midrule
He & 2 & $1s^2$ & $H{=}+1,H{=}-1$ pair, net $H{=}0$ \\
Ne & 10 & $2s^22p^6$ & oct+3 tet all filled \\
Ar & 18 & $3s^23p^6$ & $n=3$ valence closed \\
Kr & 36 & $3d^{10}4s^24p^6$ & d-shell adds 5 Hopf modes \\
Xe & 54 & $4d^{10}5s^25p^6$ & Complete period 5 \\
Rn & 86 & $4f^{14}..6p^6$ & f-shell: 7 highest modes \\
\bottomrule
\end{tabular}
\end{center}

\section{The Periodic Table as Hopf Atlas}

Each period corresponds to completing a Hopf multiplet at level $n$, 
with capacity $2n^2$. Period lengths 2, 8, 8, 18, 18, 32, 32 
follow directly. The periodic table is an atlas of Hopf multiplet 
completions, ordered by $Z$.

\section{Why Chemistry Exists}

Chemistry requires atoms with open Hopf channels: incomplete 
outer-shell multiplets that can accept or donate electrons to 
achieve topological closure. Every chemical bond is a Hopf channel 
negotiation. Covalent bonds share Hopf modes; ionic bonds transfer 
them. The driving force: topological closure.

Standard chemistry says atoms react to achieve full outer shells. 
BCT says atoms react to achieve closed Hopf multiplets. BCT is 
the reason why the first statement is true.

\section{Conclusion}

Atomic structure, the Pauli principle, the octet rule, and noble 
gas inertness all derive from BCT vacuum topology. The periodic 
table is a Hopf fibration atlas. Every noble gas is a topological 
completion. Every chemical bond is a Hopf channel negotiation.

\begin{acknowledgments}
The instinct that noble gases must be topologically complete 
structures --- not merely energetically stable --- was stated 
by Michel Cabri\'e before the calculation was done. 
The instinct was correct. The universe agreed.
\end{acknowledgments}

\begin{thebibliography}{99}
\bibitem{AppJH} M.~R.~Cabri\'e, \textit{Appendix JH: The Hopfion 
Interior and OHC}, Zenodo (2026), doi:10.5281/zenodo.18975018

\bibitem{L1} M.~R.~Cabri\'e, \textit{BCT Letter 1}, Zenodo (2026), 
doi:10.5281/zenodo.18958207

\bibitem{L71} M.~R.~Cabri\'e, \textit{BCT Letter 71: Baryogenesis}, 
Zenodo (2026).

\bibitem{White2026} H.~White, J.~Vera, A.~Sylvester, L.~Dudzinski, 
\textit{Emergent quantization from a dynamic vacuum}, 
Phys.\ Rev.\ Research \textbf{8}, 013264 (2026).
\end{thebibliography}

\end{document}
